\chapter{Trầm cảm hay Lười biếng?}

\begin{truyen}{Góc Khuất}{Mental Health vs. Laziness}
\chuhoa{H}{ọc sinh:} ``Thầy ơi, dạo này em thấy mình bị mất kết nối với bản thân, tinh thần kiệt quệ, chắc em bị trầm cảm nhẹ rồi ạ. Em xin phép không làm bài tập tuần này.''
\textit{(Student: "Teacher, lately I've been feeling disconnected from myself, mentally exhausted, I think I'm slightly depressed. I ask for permission to skip this week's homework.")}

\teacherpair{Thầy}{Vậy sao? Thầy cũng đang thấy mình bị 'mất kết nối' với bảng điểm của em. Để thầy kết nối lại bằng cách cho em 0 điểm chuyên cần tuần này nhé, coi như mình cùng nhau 'chữa lành' bằng sự kỷ luật.}
\textit{(Teacher: "Is that so? I'm also feeling 'disconnected' from your gradebook. Let me reconnect by giving you a 0 for attendance this week, so we can 'heal' together through discipline.")}
\end{truyen}

\begin{baihoc}
Đừng đánh tráo khái niệm giữa sự mệt mỏi nhất thời và các vấn đề tâm lý nghiêm trọng để làm lá chắn cho sự trì trệ.
\textit{(Do not confuse temporary fatigue with serious psychological issues to shield your stagnation.)}
\end{baihoc}

\begin{ghinhoanh}
\textbf{Reflect:}
Is it a real wound or just a lack of will?
\textit{(Đó là vết thương thật hay chỉ là sự thiếu hụt ý chí?)}
\end{ghinhoanh}

\ngancach
